\subsection{Assumption M by Fourier quadrature (any $\eps$, any $k_0$)}\label{sec:Mexact}

\begin{proposition}[exact certificate for Assumption~\ref{ass:M}]\label{prop:Mexact}\Ps{}
Let $h$ be a trapezoid centred at $c$ with plateau half-width $p$, support half-width $s$, $\rho=s-p$, $A=(s+p)/2$. Its Fourier transform is $e^{itc}H(t)$ with $H(t)=\frac4\rho\frac{\sin(At)\sin(\rho t/2)}{t^2}$, $|H|\le4/(\rho t^2)$, and by Parseval
\[
\E_\pi h(X)=\frac1\pi\int_0^\infty H(t)\,\mathrm{Re}\big[e^{itc}\overline{\varphi(t)}\big]\dd t,\qquad X=\eps\lambda,
\]
with $\varphi$ the exact characteristic function of Section~\ref{sec:tv} ($\dd L/\dd x=\tfrac k2x/E_1(x)$). Since $|\varphi|$ is non-increasing (Lemma~\ref{lem:elem}(c)), the tail beyond $T$ is at most $|\varphi(T)|\,4/(\pi\rho T)$.
\end{proposition}
Evaluated by Gauss--Legendre panels (two different panel/node choices agree to $\le2\times10^{-9}$) with the explicit tail bound $\le5\times10^{-6}$ (Table~\ref{tab:Mexact}); validated against stationary simulation ($0.15,0.05$) and the Gamma-comparator limits ($0.164,0.0444$).
\begin{table}[h]\centering\small
\begin{tabular}{@{}lccccc@{}}\toprule
$k_0$ & $\eps=0.3$ & $0.1$ & $0.03$ & $0.01$ & $0.001$\\\midrule
4 & $.1527/.0432$ & $.1570/.0425$ & $.1587/.0423$ & $.1592/.0423$ & $.1594/.0422$\\
2 & $.1528/.0463$ & $.1602/.0450$ & $.1628/.0446$ & $.1636/.0445$ & $.1639/.0444$\\
1.2 & $.1363/.0418$ & $.1418/.0403$ & $.1445/.0399$ & $.1453/.0397$ & $.1456/.0397$\\
0.6 & $.0855/.0324$ & $.1056/.0306$ & $.1088/.0301$ & $.1094/.0300$ & --\\\bottomrule
\end{tabular}
\caption{$\E_\pi h_1/\E_\pi h_2$ on $I_1=0.5\,\mathrm{ES}$, $I_2=1.8\,\mathrm{ES}$.}\label{tab:Mexact}
\end{table}
Hence $m\ge0.032$ ($k_0=0.6$), $0.040$ ($1.2$), $0.042$ ($4$), $0.044$ ($2$) for all $\eps\le0.3$ tested; Assumption~\ref{ass:M} is certified at every $\eps$ used in the regret experiments, and no condition $k_0>1$ is needed. The earlier Gamma-comparator route (Corollary~\ref{cor:M}) is superseded for this purpose. Limits: floating-point quadrature with convergence check (not interval arithmetic), per parameter value, $K=1$ only.

\subsection{What the a-priori box costs}\label{sec:boxcost}

The box $\Theta_c\ni\theta^*$ enters only through $c^\pm$ (information level $\lambda_T=mT\lambda_Q/(4c^+L)$, constant $A$ in Prop.~\ref{prop:E5}), $\kappa_1=\max(1,1/c^--1)$ (so $x/\theta_M\propto x/c^-$ in Lemma~\ref{lem:E4}), $\mu_{lo}$ (leverage condition $\bar\ell\lesssim\mu_{lo}$, i.e.\ leverage burn-in $\propto c^+/\mu_{lo}$), and $P_0$. After the bootstrap $c^\pm\to1\pm\delta$, so the leading term is box-free.

\begin{proposition}[cost of the box]\label{prop:boxcost}\Ps{}
Theorem~\ref{thm:R} holds for any compact box containing $\theta^*$ with the same leading constant; the additive constant grows at most polynomially in the box, like $c^+\max(1,1/c^-)$ in $T_1$ and $c^+/\mu_{lo}$ in the leverage burn-in, times the dependence of $P_0$ on the prior weight $\rho_0$.
\end{proposition}
\Nt{} ($K=1$, $n=3\times10^4$, corner start $(s\mu^*,\alpha^*/s)$ vs truth start): extra regret $\propto s^{2.0\text{--}2.1}$ at $\rho_0=1$ (e.g.\ $\eps=0.3$: $4.6,80,2449$ at $s=5,20,100$), reduced by a weak prior (at $s=20$: $21.1,1.7,0.35$ for $\rho_0=1,0.1,0.01$); the cost is a prior-weight effect. In the regime $\eps=c/T$ the scale of $\mu^*$ cannot be learned (Theorem~\ref{thm:V}), so there the box is a genuine requirement; for fixed $\eps$ a two-stage scheme (wide box, then a $(1+o(1))$-box around the first-stage estimate) should remove it but is not proved (\Ot{}).

\subsection{General $K$}\label{sec:generalK}

Everything in Sections~\ref{sec:regret}.1--\ref{sec:regret}.5 except the two $K=1$ devices (Assumption~\ref{ass:M} and Lemma~\ref{lem:E0}) is dimension-free: Prop.~\ref{prop:reg} needs only convexity of the polytope $\Theta_c=\{\mu\in[\mu_{lo},\mu_{hi}],\alpha_k\in[\alpha_{lo,k},\alpha_{hi,k}],\sum\alpha_k/\beta_k\le1-\eps/2\}$ and non-expansiveness of the $F$-projection (a convex QP); Prop.~\ref{prop:ratio} holds with coordinate ratios; Lemmas~\ref{lem:score}--\ref{lem:logdet}, \ref{lem:E5a}--\ref{lem:E5d}, \ref{lem:E4}, \ref{lem:C1} and Prop.~\ref{prop:E5} use only $\lambda^*\ge\mu^*$ and the time-change structure, with $d=K+1$. Lemma~\ref{lem:E0} is replaced by Lemma~\ref{lem:expmoment}, while Lemma~\ref{lem:E1}(a) holds verbatim for every $K$ with $Z_c=\sum_k(\alpha_k/\beta_k)S_k$: $\eps\Lambda_T=\mu T+nM'_T-(Z_c(T)-Z_c(0))$ (integrated key identity of Proposition~\ref{prop:keyid}).

\begin{theorem}[concentration of the information without Assumption M]\label{thm:FK}\Ps{}
Let $\lambda_*=\sup_{\lambda_0}\lambda_{\min}(N(\lambda_0))>0$ (Theorem~\ref{thm:tangent}). Pathwise $J_T\succeq\int_0^Taa\T(2/\lambda_0-\lambda^*/\lambda_0^2)\dd t$, whose entries are time integrals of the cubic polynomials $f_{ij}(S)=a_ia_j(2/\lambda_0-\lambda(S)/\lambda_0^2)$ with $\E f_{ij}=N(\lambda_0)_{ij}$. Clipping the state at $s_{\max}$ (componentwise) makes $f_{ij}$ Lipschitz in $d_\ell$ with constant $\Lambda_f(s_{\max})=O(s_{\max}^2)$; on $\{\max_{t\le T}S\le s_{\max}\}$ (probability $\ge1-T^{-\kappa}$ for $s_{\max}=O(\eps^{-1}\ln T)$, Lemma~\ref{lem:expmoment}) the clipping is inactive, and $|\E\tilde f-\E f|$ is exponentially small. The Poisson-equation/Freedman argument ($\hat u$ is $\Lambda_f/\gamma$-Lipschitz by Lemma~\ref{lem:ellcontr}, $|\Delta\mathcal M|\le\Lambda_f/\gamma$, $\langle\mathcal M\rangle\le(\Lambda_f/\gamma)^2\Lambda_T$) and a union over the $(K+1)(K+2)/2$ entries with $\tau=\lambda_*/2$ give, with probability $\ge1-3(K+1)(K+2)e^{-c_fT}-T^{-\kappa}$, $c_f\asymp\tau^2\gamma^2\eps/(\Lambda_f(s_{\max})^2\mu)$:
\[
J_T\succeq\frac{T\lambda_*}{2}\,I_{K+1}.
\]
\end{theorem}
Thus Assumption~\ref{ass:M} (a statement about $\pi$-mass of $K+1$ well-placed regions, not certifiable for $K\ge2$ by the Fourier method) is replaced by the computable number $\lambda_L>0$. Leverage: $\bar\ell(t)\le\frac{4Lc^+}{t\lambda_L\mu^*}(1+Ks_{\max}^2)=O((\ln t)^2/(\eps^2t))$.

\begin{theorem}[expected regret, general $K$]\label{thm:RK}\Ps{}
Under a known polytope $\Theta_c\ni\theta^*$ (information level $\lambda_*/2$ of Theorem~\ref{thm:FK}): for every $\delta,\eta>0$, $\E R_n\le(\frac12+O(\delta+\eta))(K+1)\ln T_n+C_K$ with $C_K$ independent of $n$; $\limsup_n\E R_n/\ln T_n\le(K+1)/2$. The proof is that of Theorem~\ref{thm:R} with the replacements above.
\end{theorem}
With the tangent level the constants are far less pessimistic than with truncation: for $\beta=(1,2)$, $\mu=1.5$, $\eps=0.4$ (0.3) the information level is $\lambda_*/2\approx0.0126$ ($0.0132$) against $\approx10^{-6}$ for the truncation device, a gain of about $10^4$ in $T_1$; they remain far from sharp (the rate $c_f$ contains $\Lambda_f(s_{\max})^{-2}$). \Nt{} ($K=2$, $\beta=(1,2)$, $\mu^*=1.5$, projection by SLSQP active in $0.2\%$ of steps, 6 replicates, $n=4\times10^4$): $\E R_n=4.41,5.29,6.14,7.28,7.99$ ($\eps=0.4$) and $4.85,5.86,6.60,7.64,8.28$ ($\eps=0.3$) at $n/16,\dots,n$; slope in $\ln T$: mean $1.3$ (both), target $(K+1)/2=1.5$; slope in $\log\det F$: $0.43$ vs $0.5$. This supports the order, not the exact constant, at these horizons.

\subsection{A minimax lower bound for any algorithm}\label{sec:lb}

An algorithm defines a predictive intensity $\hat\lambda_i(u)\ge\hat\mu_i>0$ on each gap, hence a probability law $Q^{(n)}$ on records of $n$ events; the cumulative log-loss regret equals $\KL(P_\theta^{(n)}\|Q^{(n)})=\E R_n$.

\begin{proposition}[lower bound]\label{prop:LB}\Ps{}
Let $\Theta_0\subset\mathrm{int}\,\Theta_\eps$ be a fixed compact box ($d=K+1$) and suppose an estimator $\hat\theta_n$ satisfies (U): $|\Theta_0|^{-1}\!\int_{\Theta_0}\E_\theta|\hat\theta_n-\theta|^2\dd\theta\le C_1\ln n/n$. Then for every $Q^{(n)}$
\[
\sup_{\theta\in\Theta_0}\KL(P_\theta^{(n)}\|Q^{(n)})\ \ge\ \ln|\Theta_0|+\frac d2\ln\frac n{\ln n}-\frac d2\ln\frac{2\pi eC_1}d ,
\]
and $\frac d2\ln n-O(1)$ if (U) holds with $C_1/n$.
\end{proposition}
\begin{proof}
$\sup\ge$ average over $\Theta\sim\mathrm{Unif}(\Theta_0)$; $\E_\Theta\KL(P_\Theta\|Q)=I(\Theta;X)+\KL(\bar P\|Q)\ge I(\Theta;X)\ge I(\Theta;\hat\theta_n)=h(\Theta)-h(\Theta|\hat\theta_n)$, with $h(\Theta)=\ln|\Theta_0|$ and $h(\Theta|\hat\theta_n)\le h(\Theta-\hat\theta_n)\le\frac d2\ln(2\pi e\,\E|\Theta-\hat\theta_n|^2/d)$ (Gaussian maximum entropy, AM--GM).
\end{proof}
Hypothesis (U) is delivered, for fixed $\eps$, by Theorem~\ref{thm:Cexp} and the information level ($D_n\le C(d\ln n+x)$, $\lambda_{\min}(F_n)\ge cn$ w.h.p.), with constants continuous in the parameters, hence bounded on a compact $\Theta_0$ (stated, not proved). Together with Theorem~\ref{thm:R}: $\inf_Q\sup_{\Theta_0}\E R_n=\frac d2\ln n\,(1+o(1))$, so the leading constant $d/2$ is minimax optimal. The first-order term in time units is $\frac d2\ln T+\frac12\ln\det I$, and $\det I$ stays bounded as $\eps\to0$ for $k_0>1$ (Thm.~\ref{thm:A}), consistent with the absence of a $\ln(1/\eps)$ term in the experiments; an $\eps$-explicit bound for $\eps=\eps(n)\to0$ is open (\Ot{}).


\subsection{Matrix-covariance lower bound, its $\eps$-dependence and the two-point form}\label{sec:lbmatrix}

\begin{proposition}[two-point bound in Fisher form]\label{prop:twopoint}\Pt{}
Let $\theta,\theta'$ lie in a ratio set ($c^-\le\lambda_{\theta'}/\lambda_\theta\le c^+$) and let both laws start from the same observed state. Then, with $\tilde c^{+}=\max(1,c^+)$, $\tilde c^-=\min(1,c^-)$ and $u=\theta'-\theta$,
\[
\frac{T\,u\T\bar I_Tu}{2\tilde c^{+}}\ \le\ \KL(P_\theta^T\|P_{\theta'}^T)\ \le\ \frac{T\,u\T\bar I_Tu}{2(\tilde c^{-})^2},\qquad \bar I_T=\tfrac1T\E_\theta J_T\ \ (=I(\theta)\text{ for a stationary start}).
\]
\end{proposition}
\begin{proof}
$\KL(P_\theta^T\|P_{\theta'}^T)$ is the regret of the constant predictor $\theta'$ (Lemma~\ref{lem:regid}); apply Lemma~\ref{lem:regsand} and take expectations.
\end{proof}
The bound is direction-resolved ($u\T Iu\sim\eps/(\mu(k_0-1))$ along the ridge, $\sim1/\eps$ along the top eigenvector), unlike the kernel-agnostic $\delta^2T\mu(1-n)$ bound of the literature.

\begin{proposition}[matrix lower bound]\label{prop:LBmat}\Pt{}
Let $\Theta\sim\mathrm{Unif}(\Theta_0)$ and suppose a statistic $\hat\theta$ satisfies $\E_\theta[(\hat\theta-\theta)(\hat\theta-\theta)\T]\preceq(\kappa_T/T)A^{-1}$ on $\Theta_0$, $A\succ0$. Then for every predictor $Q$
\[
\sup_{\theta\in\Theta_0}\mathrm{Reg}_T\ \ge\ I(\Theta;X)\ \ge\ \ln|\Theta_0|+\tfrac d2\ln\frac{T}{2\pi e\kappa_T}+\tfrac12\ln\det A .
\]
\end{proposition}
\begin{proof}
As Proposition~\ref{prop:LB}, with the Gaussian maximum-entropy bound $h(\Theta-\hat\theta)\le\frac12\ln\big((2\pi e)^d\det\E[(\Theta-\hat\theta)(\Theta-\hat\theta)\T]\big)$ and monotonicity of $\det$ in the Loewner order.
\end{proof}

\begin{theorem}[lower bound with explicit $\eps$-dependence]\label{thm:LBeps}\Ps{}
For online Newton, $\hat\theta=\theta_{n+1}$ satisfies the hypothesis of Proposition~\ref{prop:LBmat} with $A=N(\theta)$ (the tangent minorant of $I$) and $\kappa_T=2c^+c_Dd\ln T$ on the good event of Theorem~\ref{thm:Cexp}. Hence for $\Theta_0=[\mu(1-w),\mu(1+w)]\times[\beta(1-\tfrac32\eps),\beta(1-\tfrac12\eps)]$ ($K=1$) and $T\eps\gg\ln T$
\[
\sup_{\Theta_0}\mathrm{Reg}_T\ \ge\ \ln(T\eps)+\ln(2w\mu\beta)+\tfrac12\ln\det N_*-\ln(2\pi e\kappa_T),
\]
and the upper bound of Theorem~\ref{thm:R} with a prior matched to $\Theta_0$ is $\ln(T\eps)+\ln(2w\mu\beta)+\tfrac12\ln\det I_*+O(1)$. For $k_0>1$ both $\det I_*$ and $\det N_*$ are bounded above and below uniformly in $\eps$, so the two bounds agree in their $\eps$-dependence up to $\tfrac12\ln\ln T+O(1)$; for $k_0<1$ the upper bound carries the additional $\tfrac{1-k_0}2\ln(1/\eps)$ of $\tfrac12\ln\det I$.
\end{theorem}
